\chapter{Learning Theory — Rademacher}
%%% generated by the blueprint pipeline from TCSlib.LearningTheory.JohnsonLindenstrauss.Rademacher
\section{Overview}

This module establishes the Johnson--Lindenstrauss concentration bound for the
Rademacher (random $\pm 1/\sqrt{k}$) matrix.  It proceeds in two layers: a
distribution-agnostic sub-Gaussian concentration theorem
(\texttt{jl\_concentration\_single\_subgaussian}) is proved first, and then the
explicit Rademacher construction is shown to satisfy its hypotheses, yielding the
concrete bound $\Pr[\varepsilon\norm{x}^2 < \abs{\norm{Ax}^2 - \norm{x}^2}] \leq
2\exp(-k\varepsilon^2/8)$.

%%%%%%%%%%%%%%%%%%%%%%%%%%%%%%%%%%%%%%%%%%%%%%%%%%%%%%%%%%%%%%%%%%%%%%%%%%%%%%%
\section{Declarations}

\begin{theorem}[Distribution-agnostic JL single-vector concentration]
\lean{jl_concentration_single_subgaussian}
\label{jl_concentration_single_subgaussian}
\leanok
\uses{ProbabilityTheory.HasBernsteinMGF, jl_concentration_single_via_bernstein}
\difficulty{5}
Let $(\Omega,\mu)$ be a probability space, let $A\colon\Omega\to\bbr^{k\times d}$ be a
random matrix with $k\ge 1$ rows, and let $x\in\bbr^d$ be a nonzero vector. Write
$(Ax)_i$ for the $i$-th coordinate of the image $Ax\in\bbr^k$, and suppose that the $k$
coordinates $(Ax)_1,\dots,(Ax)_k$ are measurable and mutually independent, that each is
sub-Gaussian with variance proxy $\norm{x}^2/k$ — meaning $e^{t(Ax)_i}$ is
$\mu$-integrable and $\E_\mu\!\bigl[e^{t(Ax)_i}\bigr]\le
\exp\!\bigl((\norm{x}^2/k)\,t^2/2\bigr)$ for every $t\in\bbr$ — and that each has second
moment $\E_\mu\!\bigl[(Ax)_i^2\bigr]=\norm{x}^2/k$. Then for every $\varepsilon$ with
$0<\varepsilon<1$,
\[
\mu\!\left(\bigl\{\,\omega : \varepsilon\norm{x}^2 <
\bigl|\,\norm{Ax}^2-\norm{x}^2\,\bigr|\,\bigr\}\right)
  \;\le\; 2\exp\!\Bigl(-\tfrac{k\varepsilon^2}{8}\Bigr).
\]
\end{theorem}

\begin{definition}[Rademacher distribution on $\mathbb{R}$]
\lean{rademacherReal}
\label{rademacherReal}
\leanok
The Rademacher distribution $\mathrm{Rad}$ on $\mathbb{R}$ is the measure
$\tfrac{1}{2}\delta_{-1} + \tfrac{1}{2}\delta_{+1}$, placing equal mass $1/2$
at $-1$ and at $+1$.
\end{definition}

\begin{lemma}[Finiteness of one-half in the extended nonnegative reals]
\lean{rad_half_ne_top}
\label{rad_half_ne_top}
\leanok
\difficulty{2}
Regarded as an element of the extended nonnegative reals $[0,\infty]$, the number $1/2$
is not equal to $\infty$.
\end{lemma}

\begin{lemma}[{Rademacher samples lie in $[-1,1]$}]
\lean{rademacherReal_mem_Icc}
\label{rademacherReal_mem_Icc}
\leanok
\uses{rademacherReal}
\difficulty{3}
Let $\mathrm{Rad} = \tfrac{1}{2}\delta_{-1} + \tfrac{1}{2}\delta_{+1}$ be the Rademacher
distribution on $\mathbb{R}$. Then $\mathrm{Rad}$-almost every point $y$ satisfies $y
\in [-1, 1]$.
\end{lemma}

\begin{lemma}[Integrability under a scaled Dirac measure]
\lean{integrable_smul_dirac}
\label{integrable_smul_dirac}
\leanok
\uses{rad_half_ne_top}
\difficulty{1}
Fix a point $a \in \bbr$, and let $\delta_a$ denote the Dirac measure concentrated at
$a$. Then every function $f : \bbr \to \bbr$ is integrable with respect to the finite
measure $\tfrac{1}{2}\,\delta_a$.
\end{lemma}

\begin{lemma}[Integrability against the Rademacher measure]
\lean{integrable_rademacherReal}
\label{integrable_rademacherReal}
\leanok
\uses{integrable_smul_dirac, rademacherReal}
\difficulty{1}
Let $\mathrm{Rad} = \tfrac{1}{2}\delta_{-1} + \tfrac{1}{2}\delta_{+1}$ be the Rademacher
distribution on $\bbr$. Then every function $f \colon \bbr \to \bbr$ is integrable with
respect to $\mathrm{Rad}$.
\end{lemma}

\begin{lemma}[Mean of the Rademacher distribution]
\lean{integral_id_rademacherReal}
\label{integral_id_rademacherReal}
\leanok
\uses{integrable_smul_dirac, rademacherReal}
\difficulty{2}
Let $\mathrm{Rad} = \tfrac{1}{2}\delta_{-1} + \tfrac{1}{2}\delta_{+1}$ be the Rademacher
distribution on $\bbr$. Then the mean of the identity function against $\mathrm{Rad}$
vanishes, \[ \int_{\bbr} y \, d\mathrm{Rad}(y) = 0. \]
\end{lemma}

\begin{lemma}[Second moment of the Rademacher distribution]
\lean{integral_sq_rademacherReal}
\label{integral_sq_rademacherReal}
\leanok
\uses{integrable_smul_dirac, rademacherReal}
\difficulty{2}
Let $\mathrm{Rad} = \tfrac{1}{2}\delta_{-1} + \tfrac{1}{2}\delta_{+1}$ be the Rademacher
distribution on $\bbr$, placing equal mass $1/2$ at $-1$ and at $+1$. Then its second
moment equals one:
\[
\int_{\bbr} y^2 \, d\mathrm{Rad}(y) = 1.
\]
\end{lemma}

\begin{lemma}[Identity is sub-Gaussian with variance proxy one under Rademacher]
\lean{hasSubgaussianMGF_id_rademacherReal}
\label{hasSubgaussianMGF_id_rademacherReal}
\leanok
\uses{integral_id_rademacherReal, rademacherReal, rademacherReal_mem_Icc}
\difficulty{4}
Let $\mathrm{Rad}$ denote the Rademacher distribution on $\bbr$, placing mass $1/2$ at
each of $-1$ and $+1$, and regard the identity map as the real random variable $X$ with
law $\mathrm{Rad}$. Then $X$ is sub-Gaussian with variance proxy $1$: its moment
generating function satisfies
\[
\E\!\left[e^{tX}\right] \le e^{t^{2}/2} \qquad \text{for every } t \in \bbr.
\]
\end{lemma}

\begin{definition}[Rademacher matrix sample space]
\lean{RadΩ}
\label{RadΩ}
\leanok
The sample space for the Rademacher matrix is $\mathrm{Fin}\,k \to
\mathrm{Fin}\,d \to \mathbb{R}$, i.e.\ the type of $k \times d$ arrays of
real numbers, indexed by $(i,j) \in \mathrm{Fin}\,k \times \mathrm{Fin}\,d$.
\end{definition}

\begin{definition}[Joint Rademacher product measure]
\lean{radJointMeasure}
\label{radJointMeasure}
\leanok
\uses{RadΩ, rademacherReal}
The joint measure on $\mathrm{Rad\Omega}(k,d)$ is the product measure
$\bigotimes_{(i,j)} \mathrm{Rad}$, making all $k \cdot d$ entries i.i.d.\
Rademacher.
\end{definition}

\begin{definition}[Rademacher matrix]
\lean{radMatrix}
\label{radMatrix}
\leanok
\uses{RadΩ}
The Rademacher matrix $A : \mathrm{Rad\Omega}(k,d) \to \mathbb{R}^{k\times d}$
maps a sample $\omega$ to the matrix with entries
$A(\omega)_{ij} = \omega_{ij}/\sqrt{k}$, scaling the raw $\pm 1$ entries by
$1/\sqrt{k}$.
\end{definition}

\begin{lemma}[Measurability of the Rademacher matrix map]
\lean{measurable_radMatrix}
\label{measurable_radMatrix}
\leanok
\uses{radMatrix}
\difficulty{2}
Equip the sample space $\Omega = \{\omega : \mathrm{Fin}\,k \to \mathrm{Fin}\,d \to
\bbr\}$ of $k \times d$ real arrays with its product $\sigma$-algebra, and the space
$\bbr^{k \times d}$ of matrices likewise. Then the map $A : \Omega \to \bbr^{k\times d}$
sending $\omega$ to the matrix with entries $A(\omega)_{ij} = \omega_{ij}/\sqrt{k}$ is
measurable.
\end{lemma}

\begin{lemma}[Row marginal of the joint Rademacher measure]
\lean{radJointMeasure_row_marginal}
\label{radJointMeasure_row_marginal}
\leanok
\uses{RadΩ, radJointMeasure, rademacherReal}
\difficulty{2}
Fix natural numbers $k$ and $d$ and a row index $i \in \{0, \dots, k-1\}$, and consider
the joint measure on the set of $k \times d$ real arrays under which all $k \cdot d$
entries are independent Rademacher variables, each taking the values $-1$ and $+1$ with
probability $\tfrac12$ apiece. Then the pushforward of this measure along the map that
sends an array to its $i$-th row, an element of $\mathbb{R}^d$, equals the $d$-fold
product $\bigotimes_{j=1}^{d} \mathrm{Rad}$ of Rademacher distributions.
\end{lemma}

\begin{lemma}[Entry marginal of the joint Rademacher measure]
\lean{radJointMeasure_entry_marginal}
\label{radJointMeasure_entry_marginal}
\leanok
\uses{RadΩ, radJointMeasure, radJointMeasure_row_marginal, rademacherReal}
\difficulty{3}
Let $\mathrm{Rad} = \tfrac{1}{2}\delta_{-1} + \tfrac{1}{2}\delta_{+1}$ be the Rademacher
distribution on $\bbr$, and let $\mu$ be the product measure on the space of $k \times
d$ real arrays $\omega = (\omega_{ij})$ under which all $k \cdot d$ entries are
independent and $\mathrm{Rad}$-distributed. Then for every index $(i,j) \in
\mathrm{Fin}\,k \times \mathrm{Fin}\,d$, the pushforward of $\mu$ along the coordinate
projection $\omega \mapsto \omega_{ij}$ is exactly $\mathrm{Rad}$.
\end{lemma}

\begin{lemma}[Independence of the entries within a row]
\lean{radJointMeasure_row_iid}
\label{radJointMeasure_row_iid}
\leanok
\uses{RadΩ, radJointMeasure, radJointMeasure_entry_marginal, radJointMeasure_row_marginal, rademacherReal}
\difficulty{4}
Draw a $k \times d$ real matrix $\omega = (\omega_{ij})$ from the joint measure under
which all $k \cdot d$ entries are independent and identically distributed Rademacher
variables, each taking the values $-1$ and $+1$ with probability $1/2$. Then for every
fixed row index $i \in \{1,\dots,k\}$, the family of coordinate maps $\bigl(\omega
\mapsto \omega_{ij}\bigr)_{1 \le j \le d}$ — the $d$ entries lying in row $i$ — is
mutually independent.
\end{lemma}

\begin{lemma}[Rows of a Rademacher matrix are mutually independent]
\lean{radJointMeasure_rows_iid}
\label{radJointMeasure_rows_iid}
\leanok
\uses{RadΩ, radJointMeasure}
\difficulty{1}
Equip the space of $k \times d$ real arrays $\omega = (\omega_{ij})$ with the joint
Rademacher product measure, under which the $k \cdot d$ entries are independent and each
$\omega_{ij}$ takes the values $-1$ and $+1$ with equal probability $1/2$. For each
index $i \in \{1,\dots,k\}$, let $R_i$ be the $i$-th row, that is, the
$\mathbb{R}^d$-valued random variable $\omega \mapsto (\omega_{ij})_{j}$. Then the
family $(R_i)_{i}$ of $k$ rows is mutually independent.
\end{lemma}

\begin{lemma}[Rademacher entries are sub-Gaussian with variance proxy one]
\lean{hasSubgaussianMGF_entry}
\label{hasSubgaussianMGF_entry}
\leanok
\uses{RadΩ, hasSubgaussianMGF_id_rademacherReal, radJointMeasure, radJointMeasure_entry_marginal, rademacherReal}
\difficulty{3}
Let $\omega = (\omega_{ij})$ be a random $k \times d$ real array whose entries are
independent and identically distributed, each taking the values $-1$ and $+1$ with
probability $\tfrac{1}{2}$ apiece. Then for every index $i \in \{1,\dots,k\}$ and $j \in
\{1,\dots,d\}$, the single entry $\omega_{ij}$ is sub-Gaussian with variance proxy $1$;
that is, its moment generating function satisfies $\E\!\left[e^{t\,\omega_{ij}}\right]
\le e^{t^2/2}$ for every $t \in \bbr$.
\end{lemma}

\begin{lemma}[Scaled Rademacher entry is sub-Gaussian]
\lean{hasSubgaussianMGF_scaled_entry}
\label{hasSubgaussianMGF_scaled_entry}
\leanok
\uses{RadΩ, integral_id_rademacherReal, radJointMeasure, radJointMeasure_entry_marginal, rademacherReal, rademacherReal_mem_Icc}
\difficulty{5}
Fix natural numbers $k > 0$ and $d$, a vector $x \in \bbr^d$, a row index $i$, and a
column index $j$. Consider the joint Rademacher product measure on the space of $k
\times d$ real arrays $\omega = (\omega_{ij})$, under which every entry is an
independent $\pm 1$ Rademacher variable (equal to $-1$ or $+1$ with probability
$\tfrac12$ each). Then the scaled entry $\omega \mapsto (x_j/\sqrt{k})\,\omega_{ij}$ has
a sub-Gaussian moment generating function with variance proxy $x_j^2/k$; that is, its
moment generating function is finite and satisfies \[
\E\!\left[\exp\!\big(t\,(x_j/\sqrt{k})\,\omega_{ij}\big)\right] \le
\exp\!\left(\tfrac{1}{2}\,\frac{x_j^2}{k}\,t^2\right) \qquad\text{for all } t \in \bbr.
\]
\end{lemma}

\begin{lemma}[Rows of a Rademacher projection are sub-Gaussian]
\lean{hasSubgaussianMGF_row_proj}
\label{hasSubgaussianMGF_row_proj}
\leanok
\uses{RadΩ, hasSubgaussianMGF_scaled_entry, radJointMeasure, radJointMeasure_row_iid, radMatrix}
\difficulty{5}
Fix an integer $k \ge 1$ and a dimension $d$, let $x \in \bbr^d$ (with its Euclidean
norm), and fix a row index $i \in \{1,\dots,k\}$. Draw the entries of the Rademacher
matrix $A(\omega)_{ij} = \omega_{ij}/\sqrt{k}$ by making all $kd$ signs $\omega_{ij}$
independent Rademacher variables. Then the real random variable
\[
\omega \;\longmapsto\; (A(\omega)\,x)_i \;=\; \frac{1}{\sqrt{k}}\sum_{j} \omega_{ij}\,
x_j
\]
has a sub-Gaussian moment generating function with variance proxy $\norm{x}^2/k$: for
every $t \in \bbr$,
\[
\E\!\left[e^{\,t\,(A(\omega)x)_i}\right] \;\le\;
\exp\!\left(\frac{t^2}{2}\cdot\frac{\norm{x}^2}{k}\right).
\]
\end{lemma}

\begin{lemma}[Variance of the identity under the Rademacher distribution]
\lean{variance_id_rademacherReal}
\label{variance_id_rademacherReal}
\leanok
\uses{integral_id_rademacherReal, integral_sq_rademacherReal, rademacherReal, rademacherReal_mem_Icc}
\difficulty{3}
Let $X$ be a real random variable distributed according to the Rademacher distribution
$\mathrm{Rad} = \tfrac12\delta_{-1} + \tfrac12\delta_{+1}$, so that $X$ takes the values
$-1$ and $+1$ each with probability $1/2$. Then the variance of $X$ is
\[
\mathrm{Var}(X) = 1.
\]
\end{lemma}

\begin{lemma}[Variance of a single Rademacher entry]
\lean{variance_entry}
\label{variance_entry}
\leanok
\uses{RadΩ, radJointMeasure, radJointMeasure_entry_marginal, variance_id_rademacherReal}
\difficulty{2}
Let $\omega$ be a $k \times d$ random array whose entries are independent and
identically distributed with the Rademacher distribution $\tfrac{1}{2}\delta_{-1} +
\tfrac{1}{2}\delta_{+1}$, taking each of the values $-1$ and $+1$ with probability
$1/2$. Then for every index $(i,j)$ with $i \in \{0,\dots,k-1\}$ and $j \in
\{0,\dots,d-1\}$, the single entry $\omega_{ij}$ has variance
\[
\mathrm{Var}(\omega_{ij}) = 1.
\]
\end{lemma}

\begin{lemma}[Mean of a single Rademacher entry]
\lean{integral_entry}
\label{integral_entry}
\leanok
\uses{RadΩ, integral_id_rademacherReal, radJointMeasure, radJointMeasure_entry_marginal, rademacherReal}
\difficulty{3}
Let $\omega = (\omega_{ij})$ be a $k \times d$ array whose entries are independent
Rademacher variables, each taking the values $-1$ and $+1$ with probability $\tfrac12$
apiece. Then for every pair of indices $i \in \{1,\dots,k\}$ and $j \in \{1,\dots,d\}$,
the entry $\omega_{ij}$ has mean zero: $\E[\omega_{ij}] = 0$.
\end{lemma}

\begin{lemma}[Scaled Rademacher entry lies in $L^2$]
\lean{memLp_scaled_entry}
\label{memLp_scaled_entry}
\leanok
\uses{RadΩ, radJointMeasure, radJointMeasure_entry_marginal, rademacherReal_mem_Icc}
\difficulty{4}
Fix $k > 0$ and a vector $x \in \bbr^d$, and let $\mu$ be the joint Rademacher product
measure on the sample space of $k \times d$ arrays $\omega = (\omega_{ij})$. Then for
every pair of indices $i$ and $j$, the coordinate function $\omega \mapsto
\dfrac{x_j}{\sqrt{k}}\,\omega_{ij}$ is square-integrable with respect to $\mu$; that is,
it belongs to $L^2(\mu)$.
\end{lemma}

\begin{lemma}[Variance of a scaled Rademacher entry]
\lean{variance_scaled_entry}
\label{variance_scaled_entry}
\leanok
\uses{RadΩ, radJointMeasure, variance_entry}
\difficulty{3}
Let $k > 0$ and let $x \in \bbr^d$. Draw a $k \times d$ array $\omega$ whose entries
$\omega_{ij}$ are independent Rademacher variables, each equal to $-1$ or $+1$ with
probability $\tfrac{1}{2}$. Then for any indices $i \in \{1,\dots,k\}$ and $j \in
\{1,\dots,d\}$, the real random variable $\omega \mapsto (x_j/\sqrt{k})\,\omega_{ij}$
has variance
\[
\mathrm{Var}\!\left[\frac{x_j}{\sqrt{k}}\,\omega_{ij}\right] = \frac{x_j^2}{k}.
\]
\end{lemma}

\begin{lemma}[Variance of a row of a Rademacher projection]
\lean{variance_row_proj}
\label{variance_row_proj}
\leanok
\uses{RadΩ, memLp_scaled_entry, radJointMeasure, radJointMeasure_row_iid, radMatrix, variance_scaled_entry}
\difficulty{5}
Fix an integer $k \ge 1$, a vector $x \in \bbr^d$, and a row index $i \in
\{1,\dots,k\}$. Let $A$ be the Rademacher matrix drawn from the joint Rademacher product
measure, so its $k \cdot d$ entries $A_{ij} = \omega_{ij}/\sqrt{k}$ are independent,
each equal to $\pm 1/\sqrt{k}$ with probability $1/2$. Then the $i$-th coordinate of the
image $Ax$ has variance
\[
  \mathrm{Var}\!\bigl[(Ax)_i\bigr] = \frac{\norm{x}^2}{k}.
\]
\end{lemma}

\begin{lemma}[Row projection of the Rademacher matrix has mean zero]
\lean{integral_row_proj}
\label{integral_row_proj}
\leanok
\uses{RadΩ, integral_entry, memLp_scaled_entry, radJointMeasure, radMatrix}
\difficulty{4}
Fix integers $k > 0$ and $d \ge 0$, and let the entries $\omega_{ij}$, for $i \in
\{1,\dots,k\}$ and $j \in \{1,\dots,d\}$, be independent Rademacher random variables,
each taking the values $-1$ and $+1$ with probability $1/2$. Form the $k \times d$
random matrix $M$ with entries $M_{ij} = \omega_{ij}/\sqrt{k}$. Then for every fixed
vector $x \in \bbr^d$ and every row index $i \in \{1,\dots,k\}$, the $i$-th coordinate
of the image $Mx$ has expectation zero:
\[
\E\bigl[(Mx)_i\bigr] = 0.
\]
\end{lemma}

\begin{lemma}[Second moment of a single row of a Rademacher projection]
\lean{integral_sq_row_proj}
\label{integral_sq_row_proj}
\leanok
\uses{RadΩ, hasSubgaussianMGF_row_proj, integral_row_proj, memLp_scaled_entry, radJointMeasure, radMatrix, variance_row_proj}
\difficulty{4}
Fix integers $k > 0$ and $d \ge 0$, a vector $x \in \bbr^d$, and a row index $i \in \{1,
\dots, k\}$. Draw $\omega$ from the joint Rademacher measure on
$\mathrm{Rad\Omega}(k,d)$, form the Rademacher matrix $A(\omega) \in \bbr^{k \times d}$,
and let $(A(\omega)x)_i$ denote the $i$-th coordinate of the image of $x$ under the
linear map $A(\omega)$. Then the second moment of this coordinate is
\[
  \E\bigl[\,(A(\omega)x)_i^2\,\bigr] = \frac{\norm{x}^2}{k}.
\]
\end{lemma}

\begin{lemma}[Independence of the row projections of a Rademacher matrix]
\lean{radMatrix_proj_indep}
\label{radMatrix_proj_indep}
\leanok
\uses{RadΩ, radJointMeasure, radJointMeasure_rows_iid, radMatrix}
\difficulty{3}
Fix natural numbers $k$ and $d$, and let the entries of a $k \times d$ matrix be drawn
independently from the Rademacher distribution and each scaled by $1/\sqrt{k}$, giving a
random matrix $A(\omega) \in \bbr^{k \times d}$ whose sample $\omega$ is distributed
according to the joint Rademacher product measure. Then for every fixed vector $x \in
\bbr^{d}$, the $k$ real-valued random variables $\omega \mapsto
\bigl(A(\omega)\,x\bigr)_i$, one for each row index $i \in \{1,\dots,k\}$, are mutually
independent.
\end{lemma}

\begin{lemma}[Measurability of a row of the Rademacher projection]
\lean{radMatrix_proj_meas}
\label{radMatrix_proj_meas}
\leanok
\uses{RadΩ, radMatrix}
\difficulty{2}
Fix natural numbers $k, d$, a vector $x \in \bbr^d$, and an index $i \in \{1, \dots,
k\}$. Regard the sample space $\mathrm{Rad\Omega}(k,d)$ of $k \times d$ real arrays with
its product measurable structure, and for a sample $\omega$ let $A(\omega)$ be the
Rademacher matrix with entries $A(\omega)_{ij} = \omega_{ij}/\sqrt{k}$, acting as a
linear map $\bbr^d \to \bbr^k$. Then the map $\omega \mapsto (A(\omega)\,x)_i$, sending
each sample to the $i$-th coordinate of the image of $x$, is measurable.
\end{lemma}

\begin{theorem}[Johnson–Lindenstrauss concentration for a single Rademacher projection]
\lean{jl_concentration_single_rademacher}
\label{jl_concentration_single_rademacher}
\leanok
\uses{hasSubgaussianMGF_row_proj, integral_sq_row_proj, jl_concentration_single_subgaussian, radJointMeasure, radMatrix, radMatrix_proj_indep, radMatrix_proj_meas}
\difficulty{3}
Let $k \ge 1$ be an integer, let $x \in \bbr^d$ be a nonzero vector, and fix
$\varepsilon$ with $0 < \varepsilon < 1$. Let $A$ be the $k \times d$ matrix whose $kd$
entries are independent, each taking the values $+1/\sqrt{k}$ and $-1/\sqrt{k}$ with
probability $1/2$ apiece. Then the squared norm $\norm{Ax}^2$ concentrates about
$\norm{x}^2$, with the probability of a deviation exceeding $\varepsilon\norm{x}^2$
bounded by
\[
\Pr\!\bigl[\, \varepsilon\norm{x}^2 < \bigl|\,\norm{Ax}^2 - \norm{x}^2\,\bigr| \,\bigr]
  \;\le\; 2\exp\!\Bigl(-\frac{k\varepsilon^2}{8}\Bigr).
\]
The bound is Achlioptas's $\pm1$-entry variant of the Johnson–Lindenstrauss lemma, with
the Dasgupta–Gupta exponent.
\end{theorem}
